\section{Цветной экран: SPI, RGB, память и DMA}\label{les:lcd}
\lessonintro
  {рассчитать память кадра и обновлять экран без лишней перерисовки}
  {уроки 7, 10 и 11}
  {цветная полоса потенциометра с измерением времени рисования}
  {ST7789 240×240 с CS и логикой 3,3 В, потенциометр}

\subsection{Цена одного кадра}
OLED хранил один бит на пиксель. RGB565 занимает два байта: 5 бит красного, 6 зелёного и 5 синего.
Кадр 240×240 требует 115200 байт, а 800×480 --- 768000 байт. Двойная буферизация удваивает
эти объёмы. Дополнительно нужны память программы, задачи, сеть и служебные буферы.

По SPI передаём команды и пиксели в контроллер дисплея, обычно имеющий собственную память.
Параллельная RGB-панель требует постоянного потока пикселей и синхросигналов. ESP32-S3 имеет
аппаратную поддержку LCD, но подключение RGB-панели не выводится из SPI-примера.
Первый опыт делаем на SPI-модуле, а RGB изучаем на готовой плате с экраном.

\subsection{Подключение и библиотеки}
У модуля должны быть указанное разрешение и выведенный CS. SCK соедините с GPIO12,
MOSI с GPIO11, CS с GPIO10, DC с GPIO9, RST с GPIO8; MISO не используется.
Потенциометр: края к 3V3/GND, движок к GPIO1. OLED отключите от GPIO8/9.
Питание и подсветка зависят от модуля: LED/BL может требовать драйвера.
Не подключайте ток подсветки к сигнальному GPIO без проверки документации.

\begin{figure}[H]
\centering
\begin{circuitikz}[scale=0.9,transform shape]
  \node[blk,fill=blkcpu,minimum width=2.8cm,minimum height=5.8cm] at (0,0) {};
  \node[blk,fill=blkper,minimum width=2.8cm,minimum height=5.8cm] at (9,0) {};
  \node[font=\sffamily\small\bfseries] at (0,3.2) {ESP32-S3};
  \node[font=\sffamily\small\bfseries,align=center] at (9,3.4) {ST7789\\[-2pt]\scriptsize 240$\times$240, SPI};
  \foreach \y/\n in {2.25/3V3,1.5/GND,0.75/{\pin{12}},0/{\pin{11}},-0.75/{\pin{10}},-1.5/{\pin{9}},-2.25/{\pin{8}}}
    \node[font=\sffamily\scriptsize,anchor=east] at (1.4,\y) {\n};
  \foreach \y/\n in {2.25/VCC,1.5/GND,0.75/{SCL (SCK)},0/{SDA (MOSI)},-0.75/CS,-1.5/DC,-2.25/{RES (RST)}}
    \node[font=\sffamily\scriptsize,anchor=west] at (7.6,\y) {\n};
  \draw[thick,wred]    (1.4,2.25)  -- (7.6,2.25);
  \draw[thick,wblack]  (1.4,1.5)   -- (7.6,1.5);
  \draw[thick,wyellow] (1.4,0.75)  -- (7.6,0.75);
  \draw[thick,wgreen]  (1.4,0)     -- (7.6,0);
  \draw[thick,wblue]   (1.4,-0.75) -- (7.6,-0.75);
  \draw[thick,worange] (1.4,-1.5)  -- (7.6,-1.5);
  \draw[thick,wpurple] (1.4,-2.25) -- (7.6,-2.25);
  % потенциометр слева
  \node[font=\sffamily\scriptsize,anchor=west] at (-1.4,1.2) {3V3};
  \node[font=\sffamily\scriptsize,anchor=west] at (-1.4,0) {\pin{1}};
  \node[font=\sffamily\scriptsize,anchor=west] at (-1.4,-1.2) {GND};
  \draw[thick,wred]   (-1.4,1.2) -- (-3.6,1.2);
  \draw[thick,wblack] (-1.4,-1.2) -- (-3.6,-1.2);
  \draw[thick] (-3.6,1.2) to[R,l_=\SI{10}{k\ohm}] (-3.6,-1.2);
  \draw[thick,wblue] (-1.4,0) -- (-3.0,0);
  \draw[-latex,thick] (-3.0,0) -- (-3.6,0);
  \node[font=\sffamily\scriptsize,align=center] at (-3.4,-2.1) {потенциометр\\ (движок)};
\end{circuitikz}
\caption{Подключение ST7789 по SPI (\pin{12}, \pin{11}, \pin{10}, \pin{9}, \pin{8}) и потенциометра на \pin{1}.
Вывод MISO не используется. Подсветку (BLK) подключайте по документации модуля}
\end{figure}
Установите Adafruit GFX Library и Adafruit ST7735 and ST7789 Library.
Базовый шрифт GFX не содержит кириллицу, поэтому пример использует латинские подписи.

\codecap{Рисуем только изменившуюся часть полосы}
\codeinput{l33_tft}

Сохранённая ширина позволяет закрашивать только изменившийся прямоугольник.
Сравните это с \code{fillScreen()} на каждом шаге: уменьшатся трафик и мерцание.
При 40 МГц идеальная передача полного кадра требует не меньше
$115200\cdot8/40000000\approx23$ мс; команды и паузы увеличат время.
Нельзя требовать 60 полных кадров/с от такого канала только потому, что CPU работает на 240 МГц.

\subsection{DMA, PSRAM и время жизни кадра}
DMA разгружает CPU при передаче блоков. Пригодность внешней памяти зависит от конкретного
драйвера и режима. PSRAM не делает любой указатель подходящим для любой DMA-передачи.
Проверьте требования к размещению и выравниванию; иногда нужен внутренний передаточный буфер,
даже когда кадр хранится в PSRAM. Драйвер RGB LCD ESP-IDF предусматривает кадровые и
промежуточные bounce-буферы. Пока DMA отправляет область, её нельзя менять без синхронизации.
Наш простой GFX-пример не демонстрирует все эти режимы: он сначала даёт измеримый базовый результат.

\subsection{От рисунков к LVGL}
LVGL добавляет виджеты, стили и ввод. Ей нужны работающий драйвер экрана, функция передачи области
и обслуживание таймеров. Сначала получите правильные цвета и геометрию, затем подключайте LVGL
выбранной версии: API разных основных версий отличается. Все вызовы интерфейса выполняйте
из одной задачи либо под общей блокировкой. Задача датчика передаёт данные очередью.
Для RGB-проекта заранее составьте карту выводов: экран занимает много GPIO и может конфликтовать
с камерой, памятью и другими периферийными устройствами.

\begin{trouble}
Белый экран: питание, RESET, CS или не тот контроллер. Сдвиг: размеры и смещения матрицы.
Неверные цвета: RGB/BGR или порядок байтов. Полосы: питание, скорость либо меняющийся буфер.
\end{trouble}
\begin{summary}
Экран выбирают вместе с интерфейсом и бюджетом памяти. Частичное обновление часто полезнее
увеличения частоты; доступность PSRAM и пригодность буфера для DMA проверяют отдельно.
\end{summary}

\begin{selfcheck}
\begin{enumerate}
\item Сколько занимает кадр 320×240 RGB565?
\item Чем SPI-дисплей отличается от RGB-панели?
\item Почему PSRAM не решает все вопросы DMA?
\end{enumerate}
\end{selfcheck}
\needspace{20\baselineskip}
\begin{tasks}
\begin{enumerate}
\item \lv{1} Рассчитайте память двух разрешений RGB565.
\item \lv{1} Проверьте крайние положения полосы.
\item \lv{2} Измерьте полное и частичное обновление в одинаковых условиях.
\item \lv{2} Добавьте число и обновляйте только его область.
\item \lv{3} Организуйте задачи датчика и экрана с очередью; проверьте отзывчивость.
\item \lv{4} На готовой RGB-плате соберите два виджета LVGL и составьте бюджет памяти, кадра и буферов.
\end{enumerate}
\end{tasks}

\paragraph{Документация для продолжения}\href{https://docs.espressif.com/projects/esp-idf/en/v5.4/esp32s3/api-reference/peripherals/lcd/rgb_lcd.html}{Официальное руководство по теме}.
